\documentclass[10pt,a4paper]{amsart}
\usepackage{amsmath,amssymb,amsthm,graphicx}
\usepackage{verbatim}
\newenvironment{reference}{\comment}{\endcomment}
%\newenvironment{reference}{}{}
\newenvironment{slogan}{\comment}{\endcomment}
\newenvironment{history}{\comment}{\endcomment}
% For commutative diagrams we use Xy-pic
\usepackage[all]{xy}

% We use 2cell for 2-commutative diagrams.
\xyoption{2cell}
% Theorem environments.
%
\theoremstyle{plain}
\newtheorem{theorem}[subsection]{Theorem}
\newtheorem{proposition}[subsection]{Proposition}
\newtheorem{lemma}[subsection]{Lemma}

\theoremstyle{definition}
\newtheorem{definition}[subsection]{Definition}
\newtheorem{example}[subsection]{Example}
\newtheorem{exercise}[subsection]{Exercise}
\newtheorem{situation}[subsection]{Situation}

\theoremstyle{remark}
\newtheorem{remark}[subsection]{Remark}
\newtheorem{remarks}[subsection]{Remarks}

\numberwithin{equation}{subsection}

% Macros
%
\def\lim{\mathop{\mathrm{lim}}\nolimits}
\def\colim{\mathop{\mathrm{colim}}\nolimits}
\def\Spec{\mathop{\mathrm{Spec}}}
\def\Hom{\mathop{\mathrm{Hom}}\nolimits}
\def\Ext{\mathop{\mathrm{Ext}}\nolimits}
\def\SheafHom{\mathop{\mathcal{H}\!\mathit{om}}\nolimits}
\def\SheafExt{\mathop{\mathcal{E}\!\mathit{xt}}\nolimits}
\def\Sch{\mathit{Sch}}
\def\Mor{\mathop{\mathrm{Mor}}\nolimits}
\def\Ob{\mathop{\mathrm{Ob}}\nolimits}
\def\Sh{\mathop{\mathit{Sh}}\nolimits}
\def\NL{\mathop{N\!L}\nolimits}
\def\CH{\mathop{\mathrm{CH}}\nolimits}
\def\proetale{{pro\text{-}\acute{e}tale}}
\def\etale{{\acute{e}tale}}
\def\QCoh{\mathit{QCoh}}
\def\Ker{\mathop{\mathrm{Ker}}}
\def\Im{\mathop{\mathrm{Im}}}
\def\Coker{\mathop{\mathrm{Coker}}}
\def\Coim{\mathop{\mathrm{Coim}}}

% Boxtimes
%
\DeclareMathSymbol{\boxtimes}{\mathbin}{AMSa}{"02}

%
% Macros for moduli stacks/spaces
%
\def\QCohstack{\mathcal{QC}\!\mathit{oh}}
\def\Cohstack{\mathcal{C}\!\mathit{oh}}
\def\Spacesstack{\mathcal{S}\!\mathit{paces}}
\def\Quotfunctor{\mathrm{Quot}}
\def\Hilbfunctor{\mathrm{Hilb}}
\def\Curvesstack{\mathcal{C}\!\mathit{urves}}
\def\Polarizedstack{\mathcal{P}\!\mathit{olarized}}
\def\Complexesstack{\mathcal{C}\!\mathit{omplexes}}
% \Pic is the operator that assigns to X its picard group, usage \Pic(X)
% \Picardstack_{X/B} denotes the Picard stack of X over B
% \Picardfunctor_{X/B} denotes the Picard functor of X over B
\def\Pic{\mathop{\mathrm{Pic}}\nolimits}
\def\Picardstack{\mathcal{P}\!\mathit{ic}}
\def\Picardfunctor{\mathrm{Pic}}
\def\Deformationcategory{\mathcal{D}\!\mathit{ef}}
\usepackage{cleveref}
\crefname{originalreference}{original reference}{original references}
\Crefname{originalreference}{Original reference}{Original references}
\makeatletter
\newlabel{dualizing-lemma-local-cohomology-ss}{{dualizing-lemma-local-cohomology-ss}{source}{}{sourcepointer-1}{}}
\newlabel{dualizing-lemma-local-cohomology-ss@cref}{{[originalreference][0][]dualizing-lemma-local-cohomology-ss}{[1][0][]source}{}{}{}}
\makeatother
\title{An exact sequence comparing local cohomology along a specialization-stable subset with its nonminimal points}
\author{The Stacks Project Authors}
\date{stacks-project:local-cohomology.tex:local-snapshot-20261012; GNU FDL 1.2 or later}
\begin{document}
\maketitle
\noindent\textit{Editorial note.} All delivered mathematical text is the author's own, copied from the pinned Stacks Project chapter source sources/stacks-local-cohomology/pinned-original-source.tex; no statement or proof step was rewritten, reordered or removed.\par\medskip
\noindent\textit{Editorial note.} Preamble provenance: the Stacks Project keeps one single shared preamble for every chapter; this card takes it directly from the staged sources/stacks-local-cohomology/preamble.tex as three verbatim spans (lines 8-12, 14-18 and 154-217), so no derived file was written and nothing was added or redefined. The wrapper class and the base packages are supplied by the assembler.\par\medskip
\noindent\textit{Editorial note.} Bibliography provenance: neither the counted statement nor its proof contains a citation (this lane rescans the whole span before writing the card), so the card carries no bibliography and no entry was reconstructed.\par\medskip
\noindent\textit{Editorial note.} Selection provenance: this card was opened under the file-based control contract (CONTROL.md), with the mechanical pre-card gate run on its statement and proof interval before the card was written.\par\medskip
\noindent\textit{Allowed standard prerequisites (editorial).}\begin{itemize}
\item Local cohomology with support in a specialization-stable subset, and its description as a filtered colimit over closed subsets. The author's lemma labelled lemma-adjoint-ext (source lines 996-1004), delivered as the context span of this card. Identifies the colimit term in the displayed exact sequence.
\item The degeneration of the local cohomology spectral sequence of another chapter to short exact sequences. Invoked inside the counted proof at source line 1236 and attributed to the cited chapter's lemma (a cross-chapter pointer). Supplies the short exact sequences that the proof assembles into the displayed sequence.
\end{itemize}\medskip
\ifdefined\hypertarget\hypertarget{sourcepointer-1}{}\fi
\noindent\textit{Source reference.} Original reference dualizing-lemma-local-cohomology-ss: sources/stacks-local-cohomology/pinned-original-source.tex, source lines 428--428. Navigation pointer only: the counted text refers to the Stacks item dualizing-lemma-local-cohomology-ss, which belongs to another chapter of the project and is not part of this card. The pointer carries the target's own label as its number because this chapter's staged snapshot does not contain the other chapter's tag table; its content is not reproduced here.\par\medskip
\begin{lemma}
\label{lemma-filter-local-cohomology}
Let $A$ be a Noetherian ring. Let $T \subset \Spec(A)$
be a subset stable under specialization. Let $T' \subset T$ be
the set of nonminimal primes in $T$. Then $T'$
is a subset of $\Spec(A)$ stable under specialization
and for every $A$-module $M$ there is an exact sequence
$$
0 \to
\colim_{Z, f} H^1_f(H^{i - 1}_Z(M)) \to
H^i_{T'}(M) \to H^i_T(M) \to
\bigoplus\nolimits_{\mathfrak p \in T \setminus T'}
H^i_{\mathfrak p A_\mathfrak p}(M_\mathfrak p)
$$
where the colimit is over closed subsets $Z \subset T$
and $f \in A$ with $V(f) \cap Z \subset T'$.
\end{lemma}

\begin{lemma}
\label{lemma-adjoint-ext}
Let $A$ be a Noetherian ring. Let $T \subset \Spec(A)$
be a subset stable under specialization.
For any object $K$ of $D(A)$ we have
$$
H^i(RH^0_T(K)) = \colim_{Z \subset T\text{ closed}} H^i_Z(K)
$$
\end{lemma}

\begin{proof}
For every $Z$ and $f$ the spectral sequence of
Dualizing Complexes, Lemma \ref{dualizing-lemma-local-cohomology-ss}
degenerates to give short exact sequences
$$
0 \to H^1_f(H^{i - 1}_Z(M)) \to
H^i_{Z \cap V(f)}(M) \to H^0_f(H^i_Z(M)) \to 0
$$
We will use this without further mention below.

\medskip\noindent
Let $\xi \in H^i_T(M)$ map to zero in the direct sum.
Then we first write $\xi$ as the image of some $\xi' \in H^i_Z(M)$
for some closed subset $Z \subset T$, see Lemma \ref{lemma-adjoint-ext}.
Then $\xi'$ maps to zero in $H^i_{\mathfrak p A_\mathfrak p}(M_\mathfrak p)$
for every $\mathfrak p \in Z$, $\mathfrak p \not \in T'$.
Since there are finitely many of these primes,
we may choose $f \in A$ not contained in any of these
such that $f$ annihilates $\xi'$. Then $\xi'$
is the image of some $\xi'' \in H^i_{Z'}(M)$
where $Z' = Z \cap V(f)$. By our choice of $f$ we have
$Z' \subset T'$ and we get exactness at the penultimate spot.

\medskip\noindent
Let $\xi \in H^i_{T'}(M)$ map to zero in $H^i_T(M)$.
Choose closed subsets $Z' \subset Z$ with $Z' \subset T'$
and $Z \subset T$ such that $\xi$ comes from $\xi' \in H^i_{Z'}(M)$
and maps to zero in $H^i_Z(M)$. Then we can find $f \in A$
with $V(f) \cap Z = Z'$ and we conclude.
\end{proof}

\end{document}
